% Generated draft: requires mathematical and source review.
\documentclass[11pt]{article}
\usepackage[margin=25mm]{geometry}
\usepackage{fontspec}
\setmainfont{Latin Modern Roman}
\usepackage{amsmath,amssymb,mathtools,unicode-math}
\setmathfont{Latin Modern Math}
\usepackage{enumitem,needspace,xurl,hyperref}
\setlength{\parindent}{0pt}
\setlength{\parskip}{5pt}
\setlength{\emergencystretch}{3em}
\newcommand{\N}{\mathbb{N}}
\newcommand{\Z}{\mathbb{Z}}
\newcommand{\Q}{\mathbb{Q}}
\newcommand{\R}{\mathbb{R}}
\newcommand{\C}{\mathbb{C}}
\newcommand{\F}{\mathbb{F}}
\newcommand{\A}{\mathbb{A}}
\newcommand{\PP}{\mathbb{P}}
\newcommand{\E}{\mathbb{E}}
\newcommand{\eps}{\varepsilon}
\newcommand{\dd}{\,\mathrm{d}}
\begin{document}
\section*{UnsolvedMath ID: 800682}
\textbf{Source problem code:} catalogue entry
\subsection{Definitions and mathematical statement}
Let
\[
Z=a(z_1,z_2)\frac{\partial}{\partial z_1}+b(z_1,z_2)\frac{\partial}{\partial z_2},\qquad a,b\in\mathcal O(\C^2),
\]
be a holomorphic vector field. Write \(\operatorname{Sing}(Z)=\{a=b=0\}\) and \(M=\C^2\setminus\operatorname{Sing}(Z)\). For \(Z\not\equiv0\), let \(\mathcal L_Z\) denote the regular leaves: the maximal connected immersed Riemann surfaces in \(M\) tangent to \(Z\). Let \(\mathbb D=\{\zeta\in\C:|\zeta|<1\}\). Interpret a leaf as being covered by \(\C\), respectively \(\mathbb D\), when its universal covering Riemann surface is biholomorphic to \(\C\), respectively \(\mathbb D\). Set
\[
\mathcal P_Z=\{L\in\mathcal L_Z:\widetilde L\simeq\C\},\qquad
\mathcal H_Z=\{L\in\mathcal L_Z:\widetilde L\simeq\mathbb D\},
\]
and define their saturated unions
\[
P_Z=\bigcup_{L\in\mathcal P_Z}L,\qquad H_Z=\bigcup_{L\in\mathcal H_Z}L.
\]
Here \(\widetilde L\) denotes the universal cover. The term \emph{parabolic} is interpreted below in this uniformization sense, namely \(\widetilde L\simeq\C\); the supplied record does not explicitly define that term.

The source asks the following three questions.
\begin{enumerate}
\item[(i)] Can a holomorphic vector field on \(\C^2\) have both a positive-measure family of leaves covered by \(\C\) and a positive-measure family of leaves covered by \(\mathbb D\)? The source does not specify the measure or the space on which it is defined. If positive measure refers to the saturated unions measured by ambient four-dimensional Lebesgue measure \(\lambda_4\), the question becomes whether there exists \(Z\not\equiv0\) for which both unions are measurable and
\[
\lambda_4(P_Z)>0\qquad\text{and}\qquad\lambda_4(H_Z)>0.
\]
This is a conditional rendering of the measure qualification, not a determination of the source's intended convention.
\item[(ii)] If a non-pluripolar family of leaves is parabolic, must every leaf be parabolic? Under the interpretation that non-pluripolarity refers to the union of these leaves in \(\C^2\), this asks whether, for every \(Z\not\equiv0\),
\[
P_Z\text{ is non-pluripolar in }\C^2
\quad\Longrightarrow\quad
\mathcal P_Z=\mathcal L_Z.
\]
A set is pluripolar if locally it is contained in the \(-\infty\) locus of a plurisubharmonic function that is not identically \(-\infty\) on the relevant connected neighborhood; non-pluripolar means not pluripolar. The source does not specify whether its qualification concerns saturated unions or intersections with local transversals, so that choice remains unresolved.
\item[(iii)] Find conditions on \(Z\), or on its induced foliation, sufficient to ensure
\[
\forall L\in\mathcal L_Z,\qquad \widetilde L\simeq\mathbb D.
\]
No prescribed class of sufficient conditions is given.
\end{enumerate}
The use of regular leaves and the restriction \(Z\not\equiv0\) make explicit the usual one-dimensional foliation setting. The source does not state these conventions or explain how singular points or the identically zero vector field should be treated; those qualifications require review.
\subsection{Short English statement}
For a holomorphic vector field on \(\C^2\), ask whether positive-measure families of plane-covered and disc-covered leaves can coexist, whether a non-pluripolar family of parabolic leaves forces every leaf to be parabolic, and which conditions force every leaf to be disc-covered. The measure, the meaning of a non-pluripolar family, and the treatment of singularities require clarification.
\subsection{Sources}
\begin{itemize}
\item Local UnsolvedMath registry, permanent ID \(800682\), exact source problem code \texttt{catalogue entry}, title ``Holomorphic Foliation Theory --- Problem 9.'' The supplied record attributes the problem to Dinh and Sibony, with proposed year 2018.
\item Dinh, Sibony, \emph{Some Open Problems on Holomorphic Foliation Theory} (2018), Problem 9. Supplied source URL: \url{https://arxiv.org/abs/1807.05640}. The provenance records access on 2026-10-16 and extraction by \texttt{arxiv\_latex}; these are supplied provenance claims, not independently verified visits or source checks.
\item The supplied registry labels the problem \emph{open} and records an update on 2026-10-16. This is a supplied status label, not an independently verified mathematical or literature assessment. Supplied upstream revision: \texttt{d371ee917ac9647fbb06cd6fc529deec6f3dda82}.
\end{itemize}
\end{document}
